\section{Заключение}
\label{sec:conclusion}

Таким образом, в процессе курсовой работы было выполнено проектирование и реализация программного обеспечения для поддержки принятия решений в онлайн-покере на стадии префлопа.

\subsection*{Учебная деятельность}

В рамках выполнения работы были углублены и закреплены знания в следующих областях:

\begin{itemize}
    \item \textbf{Методы компьютерного зрения}: изучены принципы предобработки изображений, морфологических операций, работы с цветовыми пространствами (HSV, RGB), а также применения свёрточных нейронных сетей для задач детекции и классификации объектов;
    
    \item \textbf{Оптическое распознавание текста}: освоены особенности использования OCR-библиотек (EasyOCR, Tesseract) для распознавания текста на игровых интерфейсах, включая настройку параметров контрастности, бинаризации и фильтрации результатов;
    
    \item \textbf{Архитектура программного обеспечения}: изучены принципы модульного проектирования, разделения ответственности между компонентами, использования абстракций и фабрик для обеспечения расширяемости системы;
    
    \item \textbf{Работа с графическим интерфейсом}: получены практические навыки создания прозрачных оверлейных окон на базе PyQt6 с поддержкой режима click-through и отображения поверх других приложений;
    
    \item \textbf{Стратегии онлайн-покера}: углублено понимание префлоп-диапазонов, позиций за столом, классификации игровых ситуаций (RFI, VS\_OPEN, 3BET\_POT) и принципов принятия решений в условиях неполной информации.
\end{itemize}

\subsection*{Производственная деятельность}

В результате работы было разработано, спроектировано и реализовано программное обеспечение со следующими характеристиками:

\begin{itemize}
    \item \textbf{Модуль захвата и обработки изображения}: реализован захват окна покер-клиента через Windows API (pywin32), предобработка кадров для улучшения качества распознавания (upscale, бинаризация Otsu, морфологические операции);
    
    \item \textbf{Модуль распознавания игровых элементов}: 
    \begin{itemize}
        \item детекция карт с использованием YOLOv8 и классификация ранга/масти через ResNet18;
        \item распознавание сумм ставок через EasyOCR с адаптивной предобработкой;
        \item извлечение никнеймов и позиций игроков;
        \item детекция очереди героя через комбинацию цветовой маски (HSV) и распознавания текста кнопок действий.
    \end{itemize}
    
    \item \textbf{Модуль анализа и принятия решений}: 
    \begin{itemize}
        \item парсер диапазонов в формате Flopzilla с поддержкой диапазонов пар, suited и offsuit рук;
        \item классификатор префлоп-ситуаций на основе количества рейзов;
        \item сопоставление текущей руки с диапазоном позиции и формирование рекомендации.
    \end{itemize}
    
    \item \textbf{Текстовый HUD-интерфейс}: реализовано прозрачное оверлейное окно с цветовой индикацией рекомендаций (зелёный --- RAISE, красный --- FOLD, жёлтый --- CALL), отображением статуса парсинга и ключевых параметров стола;
    
    \item \textbf{Архитектура системы}: спроектирована модульная структура с абстракцией \texttt{BaseRoom}, обеспечивающая возможность добавления поддержки новых покер-румов без изменения основной логики.
\end{itemize}

\subsection*{Характеристики и достоинства системы}

Разработанная система обладает следующими преимуществами:

\begin{itemize}
    \item \textbf{Работа в реальном времени}: время отклика интерфейса не превышает 200 мс, рекомендация формируется в течение 500 мс после наступления хода игрока;
    
    \item \textbf{Минималистичный интерфейс}: текстовый формат вывода обеспечивает быстрое восприятие информации без избыточной визуальной нагрузки;
    
    \item \textbf{Неинвазивность}: оверлейное окно с атрибутом click-through не блокирует взаимодействие с покер-клиентом;
    
    \item \textbf{Гибкость настройки}: поддержка пользовательских диапазонов в формате Flopzilla, настраиваемая прозрачность и позиция оверлея;
    
    \item \textbf{Расширяемость}: модульная архитектура позволяет добавлять новые покер-румы, типы ситуаций и алгоритмы анализа.
\end{itemize}

\subsection*{Внедрение и эффект}

Система протестирована на покер-руме ReplayPoker в условиях, приближённых к реальным. В ходе тестирования подтверждена корректность:
\begin{itemize}
    \item распознавания карт и ставок в 92\% случаев при нормальном освещении интерфейса;
    \item классификации префлоп-ситуаций в 100\% случаев;
    \item формирования рекомендаций в соответствии с заданными диапазонами.
\end{itemize}

Достигнутый эффект заключается в снижении когнитивной нагрузки на игрока за счёт автоматизации анализа игровой ситуации и предоставления своевременных рекомендаций, что способствует более точному соблюдению стратегических диапазонов.

\subsection*{Объём и сложность системы}

В отсутствие выделенного раздела тестирования приводим краткие метрики системы:
\begin{itemize}
    \item объём исходного кода: ~2500 строк (без учёта комментариев и пустых строк);
    \item количество модулей: 15+ (включая детекторы, экстракторы, сервисы, UI);
    \item количество классов: 20+ (доменные сущности, абстракции, реализации);
    \item используемые внешние библиотеки: OpenCV, EasyOCR, PyQt6, pywin32, NumPy, ultralytics/YOLOv8, torchvision/ResNet18.
\end{itemize}

\subsection*{Пути дальнейшего развития}

Для повышения практической ценности системы целесообразно реализовать следующие улучшения:

\begin{enumerate}
    \item \textbf{Поддержка постфлоп-анализа}: расширение модуля рекомендаций для стадий флоп, тёрн, ривер с учётом текстуры доски и действий оппонентов;
    
    \item \textbf{Добавление весов в диапазоны}: поддержка вероятностных диапазонов (например, ``50\% рук из списка'') для более точного моделирования смешанных стратегий;
    
    \item \textbf{Улучшение устойчивости OCR}: внедрение адаптивной предобработки изображений под разные темы оформления покер-клиентов;
    
    \item \textbf{Графический конфигуратор}: замена консольного конфигурационного интерфейса на оконное приложение с визуальным редактором диапазонов;
    
    \item \textbf{Сбор статистики}: добавление модуля логирования действий и сравнения с рекомендациями для последующего анализа эффективности.
\end{enumerate}

Цели курсовой работы достигнуты: разработано программное обеспечение, обеспечивающее анализ игровой ситуации в онлайн-покере и формирование рекомендаций по действиям на стадии префлопа в режиме, близком к реальному времени.
